% Limitations. Numbers from paper/tables (scripts/build_paper.py) and PLAN.md entries B1, B4, E6, E7 and the
% 2000-task pilot. The longer threats-to-validity draft (sections/validity.tex) is in the git history.
\section{Limitations}
\label{sec:discussion}
\label{sec:validity}

\paragraph{Where \ours{} is weaker}
Two results do not favour \ours{}. On the 100-task workflows with the smallest data sizes, plain NSGA-II reaches a lower
IGD on 4--5 instances, probably because the search space is small enough for non-dominated sorting to converge. At
1000--2000 tasks, IMOEA/D reaches cheaper solutions at the cost end of the front. \ours{} keeps the higher HV there and
dominates most of the IMOEA/D front where the two overlap, but the IGD values are level and depend on which methods
define the reference front (Section~\ref{sec:results}). In a pilot study, neither a doubled budget nor a running
normalisation of the Tchebycheff rank nor extra weight vectors near the axes closed this gap. We attribute it to the
steady-state replacement per subproblem in IMOEA/D, and combining this replacement with the three-tier selection is a
natural next step.

\paragraph{Model}
The model makes simplifications that are common in offloading studies. Concurrent transfers do not share bandwidth, an
edge server divides its capacity statically among its tasks (a variant in which tasks queue gives the same ranking),
the cloud runs every task in parallel, receiving data costs the device no energy, and the problem is static and
offline. Following \citet{jin2026imoeadqn}, the main setting counts only the energy of the devices, and IMOEA/D comes
close when all three tiers are counted (Table~\ref{tab:sensitivity}). All results come from simulation.

\paragraph{Uncertainty}
The interval model assumes that every quantity stays within its bounds and that the quantities vary independently.
Under these assumptions the predicted intervals contain every realised value by construction, so the Monte Carlo study
tests the ranking of the fronts and the width of the intervals, not their validity. Correlated fluctuations, such as a
congested link that slows several transfers at once, and values outside the bounds are the cases in which the
intervals could fail, and they were not tested. The edge cost interval is wider than necessary because of the
dependency problem (Section~\ref{sec:model}). The interval widths of the workloads and data sizes in Suite B are
synthetic (5--30\%). The tasks of Suite A are crisp, as in Jin et al., so their uncertainty comes from the resources
alone. With $\gamma = 0$, \ours{} selects by midpoints alone, and the intervals affect its search only through the
crossover probabilities.

\paragraph{Baselines and protocol}
We implemented all baselines from their papers on the same encoding, operators and evaluator, with their published
settings and without retuning. Where a paper leaves a detail open, we chose the most literal reading. The original
IP-MOEA paper was not available, so its crowding follows later descriptions that use hypervolume contributions.
IMaOEACI was left out because its central indicator is defined only in an unpublished thesis, and the replay threshold
and exploration schedule of IMOEA-DQN, which its paper does not give, were set by us. We found no public code
or instance files of Jin et al., so our reimplementation of IMOEA-DQN could not be checked against their code or
against their published values on their own instances. Suite A reproduces only the sizes of their instances. The parameters of \ours{} ($T$, $p_{local}$, the number of
weight vectors and the action) were chosen on separate training instances, and IMOEA-DQN (fixed) uses the same action.
The other methods were not tuned on these instances, so the tuning favours \ours{}. The ablation limits this concern,
since even the untuned action $(1, 0, 0)$ and the other neighbourhood sizes and weight-vector counts remain better than
IMOEA/D (Section~\ref{sec:ablation}), but tuning IMOEA/D's neighbourhood size and number of weight vectors on the same
training set would be a fairer comparison. The main methods were run 30 times and the others 10 times.

\paragraph{Learned control}
The negative result of Section~\ref{sec:learned} covers three DQN designs trained for 24--48 episodes, and a different
state representation, reward or a much larger training budget might change it. The reward audit nevertheless suggests
that the obstacle lies in the problem and not in the learner. Once the population has converged, the reward of a single
step carries almost no information about which action is better, and the best fixed action leaves little to gain.
